% GENERATED FILE. Edit canonical lessons, then run npm run generate:books.
% canonical-source-hash: fnv1a64:3b710fb599d1537e
% canonical-lessons: GU-C331-upvas, GU-C331-vrat, GU-C331-prarthana, GU-C331-asirvad, GU-C331-cadhavvu, GU-R331-review-325-331

\chapter{Fasts and Vows}
\label{ch:gu-fasts-and-vows}
\hlchaptermodality{331}

\begin{chapteropening}
\textbf{By the end of this chapter:} \emph{I can say who keeps a fast and what they pray for.}

Complete the last lesson of chapter 331: I can say who keeps a fast and what they pray for.
\end{chapteropening}

\section[\texorpdfstring{\gu{ઉપવાસ} (upvās)}{ઉપવાસ (upvās)}]{\gu{ઉપવાસ} (upvās) — a day of fasting; a fast}

\label{lesson:GU-C331-upvas}



\begin{quote}
Before the new one: say the Gujarati for a hymn, then the Gujarati for a devotee.
\end{quote}

\subsection*{You'll want to know: \gu{ઉપવાસ}}
\textbf{\gu{ઉપવાસ}} — \emph{upvās} — \textquotedblleft{}a day of fasting; a fast\textquotedblright{}.

\begin{quote}
\gu{દાદી} \gu{દરેક} \gu{સોમવારે} \gu{ઉપવાસ} \gu{કરે} \gu{છે}.
\end{quote}

\begin{quote}
\emph{dādī darek somvāre upvās kare che.}
\end{quote}

\begin{quote}
My grandmother fasts every Monday.
\end{quote}

A masculine noun.

\begin{cousinweb}
\textbf{\gu{ઉપવાસ}} (\emph{upvās}) is a tatsama, Sanskrit \emph{upavāsa}, \textquotedblleft{}fasting\textquotedblright{} (CDIAL 2244). It is kept with \textbf{\gu{કરવું}}: \textbf{\gu{ઉપવાસ} \gu{કરે} \gu{છે}}, fasts.
\end{cousinweb}

\begin{grammarlens}[title={What you've built}]
One more word for telling what happened.
\end{grammarlens}

\subsection*{Your turn}
\begin{itemize}
\raggedright
  \item \emph{Say it:} \emph{upvās}
  \item \emph{Say it:} \emph{upvās}, once more
  \item \emph{Say it:} \emph{dādī darek somvāre upvās kare che.} — My grandmother fasts every Monday.
  \item \emph{Recall:} say the Gujarati for a hymn, then the Gujarati for a devotee, then say \emph{upvās} again
\end{itemize}

\begin{tcolorbox}[breakable,colback=teal!4,colframe=teal!35!black,title={Before you move on}]
What does \gu{ઉપવાસ} mean? (\textquotedblleft{}A day of fasting; a fast\textquotedblright{}.) Say it once more.
\end{tcolorbox}
\section[\texorpdfstring{\gu{વ્રત} (vrat)}{વ્રત (vrat)}]{\gu{વ્રત} (vrat) — a vow (of worship)}

\label{lesson:GU-C331-vrat}



\begin{quote}
Before the new one: say the Gujarati for a devotee, then the Gujarati for a day of fasting.
\end{quote}

\subsection*{You'll want to know: \gu{વ્રત}}
\textbf{\gu{વ્રત}} — \emph{vrat} — \textquotedblleft{}a vow (of worship)\textquotedblright{}.

\begin{quote}
\gu{માતાએ} \gu{દીકરા} \gu{માટે} \gu{વ્રત} \gu{રાખ્યું}.
\end{quote}

\begin{quote}
\emph{mātāe dīkrā māṭe vrat rākhyũ.}
\end{quote}

\begin{quote}
The mother kept a vow for her son.
\end{quote}

A neuter noun.

\begin{cousinweb}
\textbf{\gu{વ્રત}} (\emph{vrat}) is a tatsama, Sanskrit \emph{vrata}, \textquotedblleft{}a vow, a religious observance\textquotedblright{} (Monier-Williams). It is kept, \textbf{\gu{વ્રત} \gu{રાખ્યું}}, and often means fasting on set days.
\end{cousinweb}

\begin{grammarlens}[title={What you've built}]
One more word for telling what happened.
\end{grammarlens}

\subsection*{Your turn}
\begin{itemize}
\raggedright
  \item \emph{Say it:} \emph{vrat}
  \item \emph{Say it:} \emph{vrat}, once more
  \item \emph{Say it:} \emph{mātāe dīkrā māṭe vrat rākhyũ.} — The mother kept a vow for her son.
  \item \emph{Recall:} say the Gujarati for a devotee, then the Gujarati for a day of fasting, then say \emph{vrat} again
\end{itemize}

\begin{tcolorbox}[breakable,colback=teal!4,colframe=teal!35!black,title={Before you move on}]
What does \gu{વ્રત} mean? (\textquotedblleft{}A vow (of worship)\textquotedblright{}.) Say it once more.
\end{tcolorbox}
\section[\texorpdfstring{\gu{પ્રાર્થના} (prārthanā)}{પ્રાર્થના (prārthanā)}]{\gu{પ્રાર્થના} (prārthanā) — a prayer}

\label{lesson:GU-C331-prarthana}



\begin{quote}
Before the new one: say the Gujarati for a day of fasting, then the Gujarati for a vow.
\end{quote}

\subsection*{You'll want to know: \gu{પ્રાર્થના}}
\textbf{\gu{પ્રાર્થના}} — \emph{prārthanā} — \textquotedblleft{}a prayer\textquotedblright{}.

\begin{quote}
\gu{શાળામાં} \gu{રોજ} \gu{સવારે} \gu{પ્રાર્થના} \gu{થાય} \gu{છે}.
\end{quote}

\begin{quote}
\emph{śāḷāmā̃ roj savāre prārthanā thāy che.}
\end{quote}

\begin{quote}
There is a prayer at school every morning.
\end{quote}

A feminine noun.

\begin{cousinweb}
\textbf{\gu{પ્રાર્થના}} (\emph{prārthanā}) is a tatsama, Sanskrit \emph{prārthanā}, \textquotedblleft{}a wish, a request, a prayer\textquotedblright{} (Monier-Williams). It is feminine: \textbf{\gu{પ્રાર્થના} \gu{કરી}}, prayed.
\end{cousinweb}

\begin{grammarlens}[title={What you've built}]
One more word for telling what happened.
\end{grammarlens}

\subsection*{Your turn}
\begin{itemize}
\raggedright
  \item \emph{Say it:} \emph{prārthanā}
  \item \emph{Say it:} \emph{prārthanā}, once more
  \item \emph{Say it:} \emph{śāḷāmā̃ roj savāre prārthanā thāy che.} — There is a prayer at school every morning.
  \item \emph{Recall:} say the Gujarati for a day of fasting, then the Gujarati for a vow, then say \emph{prārthanā} again
\end{itemize}

\begin{tcolorbox}[breakable,colback=teal!4,colframe=teal!35!black,title={Before you move on}]
What does \gu{પ્રાર્થના} mean? (\textquotedblleft{}A prayer\textquotedblright{}.) Say it once more.
\end{tcolorbox}
\section[\texorpdfstring{\gu{આશીર્વાદ} (āśīrvād)}{આશીર્વાદ (āśīrvād)}]{\gu{આશીર્વાદ} (āśīrvād) — a blessing}

\label{lesson:GU-C331-asirvad}



\begin{quote}
Before the new one: say the Gujarati for a vow, then the Gujarati for a prayer.
\end{quote}

\subsection*{You'll want to know: \gu{આશીર્વાદ}}
\textbf{\gu{આશીર્વાદ}} — \emph{āśīrvād} — \textquotedblleft{}a blessing\textquotedblright{}.

\begin{quote}
\gu{મારા} \gu{દાદાએ} \gu{મને} \gu{આશીર્વાદ} \gu{આપ્યા}.
\end{quote}

\begin{quote}
\emph{mārā dādāe mane āśīrvād āpyā.}
\end{quote}

\begin{quote}
My grandfather gave me his blessing.
\end{quote}

A masculine noun.

\begin{cousinweb}
\textbf{\gu{આશીર્વાદ}} (\emph{āśīrvād}) is a tatsama, Sanskrit \emph{āśīrvāda}, \textquotedblleft{}a benediction, a spoken blessing\textquotedblright{} (Monier-Williams). Elders give blessings in the plural: \textbf{\gu{આશીર્વાદ} \gu{આપ્યા}}.
\end{cousinweb}

\begin{grammarlens}[title={What you've built}]
One more word for telling what happened.
\end{grammarlens}

\subsection*{Your turn}
\begin{itemize}
\raggedright
  \item \emph{Say it:} \emph{āśīrvād}
  \item \emph{Say it:} \emph{āśīrvād}, once more
  \item \emph{Say it:} \emph{mārā dādāe mane āśīrvād āpyā.} — My grandfather gave me his blessing.
  \item \emph{Recall:} say the Gujarati for a vow, then the Gujarati for a prayer, then say \emph{āśīrvād} again
\end{itemize}

\begin{tcolorbox}[breakable,colback=teal!4,colframe=teal!35!black,title={Before you move on}]
What does \gu{આશીર્વાદ} mean? (\textquotedblleft{}A blessing\textquotedblright{}.) Say it once more.
\end{tcolorbox}
\section[\texorpdfstring{\gu{ચઢાવવું} (caḍhāvvũ)}{ચઢાવવું (caḍhāvvũ)}]{\gu{ચઢાવવું} (caḍhāvvũ) — to offer (to a god); to raise}

\label{lesson:GU-C331-cadhavvu}



\begin{quote}
Before the new one: say the Gujarati for a prayer, then the Gujarati for a blessing.
\end{quote}

\subsection*{You'll want to know: \gu{ચઢાવવું}}
\textbf{\gu{ચઢાવવું}} — \emph{caḍhāvvũ} — \textquotedblleft{}to offer (to a god); to raise\textquotedblright{}.

\begin{quote}
\gu{ભક્તે} \gu{મંદિરમાં} \gu{ફૂલ} \gu{ચઢાવ્યાં}.
\end{quote}

\begin{quote}
\emph{bhakte mandirmā̃ phūl caḍhāvyā̃.}
\end{quote}

\begin{quote}
The devotee offered flowers at the temple.
\end{quote}

Past, with \gu{તેણે} (he or she …): \gu{તેણે} \gu{ફૂલ} \gu{ચઢાવ્યાં}, \emph{teṇe phūl caḍhāvyā̃}.

\begin{cousinweb}
\textbf{\gu{ચઢાવવું}} (\emph{caḍhāvvũ}) is the causative of \textbf{\gu{ચઢવું}}, to climb, from \emph{*caḍhati}, \textquotedblleft{}rises\textquotedblright{}, a form CDIAL reconstructs (CDIAL 4578): to make something go up, and so to offer it up to a god.
\end{cousinweb}

\begin{grammarlens}[title={What you've built}]
One more word for telling what happened.
\end{grammarlens}

\subsection*{Your turn}
\begin{itemize}
\raggedright
  \item \emph{Say it:} \emph{caḍhāvvũ}
  \item \emph{Say it:} \emph{caḍhāvvũ}, once more
  \item \emph{Say it:} \emph{bhakte mandirmā̃ phūl caḍhāvyā̃.} — The devotee offered flowers at the temple.
  \item \emph{Recall:} say the Gujarati for a prayer, then the Gujarati for a blessing, then say \emph{caḍhāvvũ} again
\end{itemize}

\begin{tcolorbox}[breakable,colback=teal!4,colframe=teal!35!black,title={Before you move on}]
What does \gu{ચઢાવવું} mean? (\textquotedblleft{}To offer (to a god); to raise\textquotedblright{}.) Say it once more.
\end{tcolorbox}
\section[(dialogue)]{Looking back: from The Fight for Freedom to Fasts and Vows}

\label{lesson:GU-R331-review-325-331}



\begin{quote}
Say the Gujarati for freedom and a slave.
\end{quote}

\subsection*{Your turn}
Say these aloud:

\begin{itemize}
\raggedright
  \item freedom, a slave, an Englishman, a campaign, a revolt
  \item a flag, to hoist, peace, a border, to found
  \item a celebration, to celebrate, a tradition, a custom, a holy day
  \item worship, an aarti, a blessed offering, to light, incense
  \item a rangoli, to adorn, to sparkle, a garland, a procession
  \item a pavilion, to whirl round, to swing, a hymn, a devotee
  \item a day of fasting, a vow, a prayer, a blessing, to offer (to a god)
\end{itemize}

\begin{tcolorbox}[breakable,colback=teal!4,colframe=teal!35!black,title={Before you move on}]
Freedom? (\textbf{\gu{આઝાદી}}, \emph{ājhādī}.) A slave? (\textbf{\gu{ગુલામ}}, \emph{gulām}.) An Englishman? (\textbf{\gu{અંગ્રેજ}}, \emph{angrej}.) A campaign? (\textbf{\gu{આંદોલન}}, \emph{ā̃dolan}.) A revolt? (\textbf{\gu{બળવો}}, \emph{baḷvo}.) A flag? (\textbf{\gu{ઝંડો}}, \emph{jhanḍo}.) To hoist? (\textbf{\gu{ફરકાવવું}}, \emph{pharkāvvũ}.) Peace? (\textbf{\gu{શાંતિ}}, \emph{śānti}.) A border? (\textbf{\gu{સરહદ}}, \emph{sarhad}.) To found? (\textbf{\gu{સ્થાપવું}}, \emph{sthāpvũ}.) A celebration? (\textbf{\gu{ઉત્સવ}}, \emph{utsav}.) To celebrate? (\textbf{\gu{ઉજવવું}}, \emph{ujavvũ}.) A tradition? (\textbf{\gu{પરંપરા}}, \emph{paramparā}.) A custom? (\textbf{\gu{રિવાજ}}, \emph{rivāj}.) A holy day? (\textbf{\gu{પર્વ}}, \emph{parva}.) Worship? (\textbf{\gu{પૂજા}}, \emph{pūjā}.) An aarti? (\textbf{\gu{આરતી}}, \emph{ārtī}.) A blessed offering? (\textbf{\gu{પ્રસાદ}}, \emph{prasād}.) To light? (\textbf{\gu{પ્રગટાવવું}}, \emph{pragṭāvvũ}.) Incense? (\textbf{\gu{ધૂપ}}, \emph{dhūp}.) A rangoli? (\textbf{\gu{રંગોળી}}, \emph{rangoḷī}.) To adorn? (\textbf{\gu{શણગારવું}}, \emph{śaṇgārvũ}.) To sparkle? (\textbf{\gu{ઝગમગવું}}, \emph{jhagmagvũ}.) A garland? (\textbf{\gu{માળા}}, \emph{māḷā}.) A procession? (\textbf{\gu{શોભાયાત્રા}}, \emph{śobhāyātrā}.) A pavilion? (\textbf{\gu{મંડપ}}, \emph{manḍap}.) To whirl round? (\textbf{\gu{ઘૂમવું}}, \emph{ghūmvũ}.) To swing? (\textbf{\gu{ઝૂલવું}}, \emph{jhūlvũ}.) A hymn? (\textbf{\gu{ભજન}}, \emph{bhajan}.) A devotee? (\textbf{\gu{ભક્ત}}, \emph{bhakt}.) A day of fasting? (\textbf{\gu{ઉપવાસ}}, \emph{upvās}.) A vow? (\textbf{\gu{વ્રત}}, \emph{vrat}.) A prayer? (\textbf{\gu{પ્રાર્થના}}, \emph{prārthanā}.) A blessing? (\textbf{\gu{આશીર્વાદ}}, \emph{āśīrvād}.) To offer (to a god)? (\textbf{\gu{ચઢાવવું}}, \emph{caḍhāvvũ}.)
\end{tcolorbox}
